% Generated draft: requires mathematical and source review.
\documentclass[11pt]{article}
\usepackage[margin=25mm]{geometry}
\usepackage{fontspec}
\setmainfont{Latin Modern Roman}
\usepackage{amsmath,amssymb,mathtools,unicode-math}
\setmathfont{Latin Modern Math}
\usepackage{enumitem,needspace,xurl,hyperref}
\setlength{\parindent}{0pt}
\setlength{\parskip}{5pt}
\setlength{\emergencystretch}{3em}
\newcommand{\N}{\mathbb{N}}
\newcommand{\Z}{\mathbb{Z}}
\newcommand{\Q}{\mathbb{Q}}
\newcommand{\R}{\mathbb{R}}
\newcommand{\C}{\mathbb{C}}
\newcommand{\F}{\mathbb{F}}
\newcommand{\A}{\mathbb{A}}
\newcommand{\PP}{\mathbb{P}}
\newcommand{\E}{\mathbb{E}}
\newcommand{\eps}{\varepsilon}
\newcommand{\dd}{\,\mathrm{d}}
\begin{document}
\section*{UnsolvedMath ID: 30001757}
\textbf{Source problem code:} catalogue entry
\subsection{Definitions and mathematical statement}
To interpret the question, one must specify a ground field \(k\) and a differential graded algebra \((A,d_A)\) whose underlying algebra is denoted in the supplied record by \(\bigwedge(x_1,\ldots,x_r)\). The record does not specify the field, the degrees of the generators, the meaning of this wedge notation, or the differential \(d_A\).

For a specified differential graded algebra \((A,d_A)\), a differential graded \(A\)-module \((M,D)\) has \(D^2=0\) and, in cohomological grading, satisfies
\[
D(am)=d_A(a)m+(-1)^{|a|}aD(m)
\]
for homogeneous \(a\in A\). Such a module is semifree if it admits an exhaustive increasing filtration by differential graded submodules
\[
0=F_{-1}M\subseteq F_0M\subseteq F_1M\subseteq\cdots,
\qquad M=\bigcup_{j\geq0}F_jM,
\]
whose successive quotients are direct sums of degree shifts of \((A,d_A)\). In the displayed presentation \(M=A\otimes_k V\), \(V\) denotes the graded vector space of free generators. Dimensions below mean total dimensions over \(k\), so that
\[
\dim_k H^*(M,D)=\sum_n\dim_k H^n(M,D),
\qquad
\dim_k V=\sum_n\dim_k V^n.
\]
These are standard interpretive conventions; the source record does not supply the grading or filtration conventions.

The supplied question asks whether, for every number \(r\) of generators and every semifree module of the displayed form,
\[
(M,D)=\left(\bigwedge(x_1,\ldots,x_r)\otimes_k V,D\right),
\qquad
\dim_k H^*(M,D)<\infty,
\]
one must have
\[
\dim_k V\geq 2^r.
\]
This is a transcription of the intended implication, subject to unresolved hypotheses. In particular, the record does not say whether the wedge denotes an exterior algebra or a free graded commutative algebra, whether the generators have prescribed parity or degree, or whether the module must satisfy additional boundedness or minimality conditions. Finite-dimensional cohomology permits zero cohomology, and no nonvanishing condition on \(H^*(M,D)\) or exclusion of the zero module is stated. These omissions must be resolved before treating the displayed implication as a fully specified conjecture. No additional hypotheses are imposed here.
\subsection{Short English statement}
The supplied problem asks whether a semifree module over an algebra written as \(\bigwedge(x_1,\ldots,x_r)\), with finite total cohomology, must have at least \(2^r\) free generators. The field, algebra convention, generator degrees, and possible nonvanishing hypotheses remain unspecified. The supplied curation describes this as an algebraic form of the toral-rank conjecture and labels it open; those claims have not been independently verified.
\subsection{Sources}
\begin{itemize}
\item Local registry: the supplied UnsolvedMath record, permanent ID \(30001757\), exact source problem code \texttt{catalogue entry}. The project identifies \texttt{data/unsolvedmath/problems.json} as its main registry and \texttt{catalog/problems.json} as an additional registry. Their contents were not independently inspected for this draft.
\item Supplied report citation: \emph{Arbeitsgemeinschaft: Rational Homotopy Theory in Mathematics and Physics} (2011), Oberwolfach Reports, DOI 10.4171/owr/2011/18, pp.\ 993--1045. \url{https://doi.org/10.4171/owr/2011/18}. The record attributes the problem to workshop open-problem discussions and identifies no individual proposer.
\item Supplied literature source: \emph{Rational Homotopy Theory in Mathematics and Physics, OWR 18/2011}. \url{https://doi.org/10.4171/owr/2011/18}.
\item Supplied literature source, with its supplied descriptive label: Bettiol--Mendes--T\"oben, \emph{bounds on sectional curvature}. \url{https://arxiv.org/abs/1407.6201}. Its relevance to the displayed module inequality is not explained in the record.
\item Literature assessment checked at \texttt{2026-08-21}: the supplied curation labels the problem open and describes the displayed inequality as an algebraic form of the toral-rank conjecture whose general case remains unresolved. This is supplied dated curation, not an independently verified literature assessment. The background explicitly requests independent verification of time-sensitive status and cited claims.
\item The supplied record has statement status \texttt{corrected\_verified} and reports that the mathematical content was checked against the official Oberwolfach report, with notation, context, and question form normalized without changing the source-level claim. This verification claim is preserved as registry attribution and was not independently checked here.
\end{itemize}
\end{document}
